% !TEX program = lualatex
% ECON 201, Worksheet 8: Other Elasticities. Compile with LuaLaTeX so the
% PDF is tagged (see syllabus/Econ201-Syllabus.tex for the same setup).
% Build from the course root:
%   latexmk -lualatex -cd worksheets/worksheet08.tex && latexmk -c -cd worksheets/worksheet08.tex
% Every number here is computed and printed by slides/figures/make_lecture8.py
% so the worksheet and the deck quote the same figures (Maya's month).
% Lecture 8 is a quiz day, so this worksheet is two short activities with no
% back-side problem.
\DocumentMetadata{
  pdfstandard=UA-2,
  pdfversion=2.0,
  lang=en-US,
  tagging=on
}
\documentclass[11pt]{article}

\usepackage[letterpaper, margin=0.75in, top=0.5in, bottom=0.45in]{geometry}
\usepackage{fontspec}
\setmainfont{Fira Sans}
\usepackage{amsmath}
\usepackage{amssymb}
\usepackage{graphicx}
\usepackage{booktabs}
\usepackage{array}
\usepackage{xcolor}
\usepackage{needspace}
\usepackage[most]{tcolorbox}
\definecolor{cream}{HTML}{FDF3EA}
\definecolor{accent}{HTML}{BF5700}
\usepackage{titlesec}
\titleformat{\section}{\large\bfseries\color{accent}}{}{0pt}{}
\titlespacing*{\section}{0pt}{4pt}{2pt}
\usepackage{hyperref}
\hypersetup{pdftitle={ECON 201 Worksheet: Other Elasticities}, pdfauthor={Div Bhagia}, hidelinks}
\pagestyle{empty}
\setlength{\parindent}{0pt}
\setlength{\parskip}{2pt}

% Ruled answer space: n lines. Plain TeX loop; pgffor's \foreach breaks the
% enumerate counter when used inside an item.
\newcount\answerlinecount
\newcommand{\answerlines}[1]{%
  \par\vspace{1pt}\answerlinecount=#1
  \loop\ifnum\answerlinecount>0
    \rule{\linewidth}{0.3pt}\par\vspace{5pt}%
    \advance\answerlinecount by -1
  \repeat
  \vspace{-5pt}}

% Ruled grids rather than booktabs tables: students write in the cells, so
% every cell gets visible borders and the blank rows get extra height. One
% column layout for both activities, so the two tables line up.
\newcolumntype{C}{>{\centering\arraybackslash}p{2.6cm}}
\newcommand{\blankrow}{\rule{0pt}{1.8em}}

\begin{document}

{\LARGE\bfseries\color{accent} Worksheet: Other Elasticities}\hfill{\large ECON 201}

% Terms sit on the cream panel, where the brand orange falls below 4.5:1;
% the darker orange keeps WCAG AA contrast (see CLAUDE.md).
\definecolor{accentdark}{HTML}{8F4100}
\newcommand{\term}[1]{\textcolor{accentdark}{\textbf{#1}}}

% Definitions panel list: flush left, compact. Built on the generic list
% environment so the bullet sits at the panel's left edge; enumitem is avoided
% for the tagging reason noted in worksheet06.
\newenvironment{deflist}
  {\begin{list}{\textbullet}{%
     \setlength{\leftmargin}{1.1em}\setlength{\labelwidth}{0.7em}%
     \setlength{\labelsep}{0.4em}\setlength{\itemindent}{0pt}%
     \setlength{\listparindent}{0pt}\setlength{\topsep}{2pt}%
     \setlength{\partopsep}{0pt}\setlength{\parsep}{0pt}%
     \setlength{\itemsep}{2pt}}}
  {\end{list}}


\begin{tcolorbox}[colback=cream, colframe=accent, boxrule=0.6pt, arc=2pt,
  left=6pt, right=6pt, top=4pt, bottom=4pt, title=Definitions,
  fonttitle=\bfseries, coltitle=white]
\begin{deflist}
\item \term{Income elasticity of demand}: the percentage change in demand that
  would occur in response to a 1\% increase in income,
  \[
    \varepsilon_{\text{income}}
    = \frac{\%\text{ change in demand}}{\%\text{ change in income}}
  \]
  \begin{deflist}
  \item \term{Normal good}: $\varepsilon_{\text{income}} > 0$. People buy more
    of it as their income rises. There are two types:
    \begin{deflist}
    \item \term{Necessity}: $0 < \varepsilon_{\text{income}} < 1$. Demand grows
      more slowly than income.
    \item \term{Luxury}: $\varepsilon_{\text{income}} > 1$. Demand grows faster
      than income.
    \end{deflist}
  \item \term{Inferior good}: $\varepsilon_{\text{income}} < 0$. People buy less
    of it as their income rises.
  \end{deflist}
\item \term{Cross-price elasticity of demand} for good A: the percentage change
  in demand for A that would occur in response to a 1\% increase in the price
  of good B,
  \[
    \varepsilon_{\text{cross}}
    = \frac{\%\text{ change in demand for A}}{\%\text{ change in the price of B}}
  \]
  \begin{deflist}
  \item \term{Substitutes}: $\varepsilon_{\text{cross}} > 0$. A rise in the
    price of B raises demand for A.
  \item \term{Complements}: $\varepsilon_{\text{cross}} < 0$. A rise in the
    price of B lowers demand for A.
  \end{deflist}
\end{deflist}
\end{tcolorbox}

\needspace{12\baselineskip}\section{1. Income elasticity}

Maya works part time, and her monthly income rises from \$1,000 to \$1,200.
Her purchases for the month change as below.

\begin{center}
\renewcommand{\arraystretch}{1.4}
\begin{tabular}{|>{\raggedright\arraybackslash}p{3.1cm}|C|C|C|C|C|}
\hline
 & At \$1,000 & At \$1,200 & \% change in demand & Income elasticity & Normal or inferior? \\
\hline
Packs of ramen & 40 & 32 & \rule{0pt}{1.8em} & & \\
\hline
Gallons of gas & 20 & 22 & \rule{0pt}{1.8em} & & \\
\hline
Steak dinners & 2 & 3 & \rule{0pt}{1.8em} & & \\
\hline
\end{tabular}
\end{center}

Which of the normal goods is a necessity, and which is a luxury?
\answerlines{1}

\needspace{12\baselineskip}\section{2. Cross-price elasticity}

The campus cafe raises the price of a coffee from \$2.50 to \$3.00, and Maya's
purchases for the month change as below.

\begin{center}
\renewcommand{\arraystretch}{1.4}
\begin{tabular}{|>{\raggedright\arraybackslash}p{3.1cm}|C|C|C|C|C|}
\hline
 & At \$2.50 & At \$3.00 & \% change in demand & Cross-price elasticity & Substitutes or complements? \\
\hline
Coffee across the street & 6 & 9 & \rule{0pt}{1.8em} & & \\
\hline
Donuts at the campus cafe & 8 & 6 & \rule{0pt}{1.8em} & & \\
\hline
Packs of ramen & 32 & 32 & \rule{0pt}{1.8em} & & \\
\hline
\end{tabular}
\end{center}

What does ramen's cross-price elasticity tell you?
\answerlines{1}

\end{document}
